% =====================================================================
%  1bat-ud4-15.tex  —  SESSIÓ 15: Taller de preparació de la prova
%  (sense preàmbul: s'inclou des de main.tex)
% =====================================================================
\horatitol{Sessió 15 — Taller de preparació de la prova}%
{Conèixer l'estructura de la prova i la rúbrica, assajar un simulacre per blocs i fer un pla d'estudi personal.}

\subsection{Idea clau}

Demà tindrem el repàs ben fresc; avui ens hi preparem \textbf{sense sorpreses}. La
prova de la sessió següent té \textbf{tres blocs}, un per cada criteri que hem
treballat (C1, C2 i C3). Saber \emph{exactament} què es demana i com es corregeix fa
que l'examen deixi de ser imprevisible. Després assajarem amb un \textbf{simulacre}.

\begin{idea}
\textbf{Estructura de la prova} (un bloc per criteri):
\begin{enumerate}[leftmargin=1.6em, itemsep=2pt]
  \item \textbf{Bloc 1 (C1) — Paràmetres estadístics:} calcular i interpretar mitjana,
        mediana, moda i desviació típica d'un conjunt de dades.
  \item \textbf{Bloc 2 (C2) — Lectura i lectura crítica de gràfics:} llegir una
        piràmide, un histograma o un percentatge i detectar-ne els paranys.
  \item \textbf{Bloc 3 (C3) — Enquestes:} avaluar o redactar una pregunta neutra i
        tabular un petit conjunt de respostes.
\end{enumerate}
\textbf{Es valora}, a més del resultat: la \textbf{interpretació} dins del context
social, la \textbf{mirada crítica} i la \textbf{neutralitat} de les preguntes.
\end{idea}

\subsection{Rúbrica simplificada}

\noindent Així es corregeix cada bloc (nivells d'assoliment):
\begin{center}
\renewcommand{\arraystretch}{1.3}
\begin{tabular}{p{2.4cm} p{10.2cm}}
\toprule
\textbf{Nivell} & \textbf{Què demostra} \\
\midrule
\textbf{AE} (excel·lent) & Calcula sense errors i \textbf{interpreta} el resultat dins
del context social. \\
\textbf{AN} (notable) & Calcula correctament amb algun error menor; la interpretació és
essencialment bona. \\
\textbf{AS} (suficient) & Fa el procediment bàsic amb alguns errors; la interpretació
és incompleta. \\
\textbf{NA} (no assolit) & No identifica el procediment o comet errors que
n'invaliden el resultat. \\
\bottomrule
\end{tabular}
\end{center}

\subsection{Simulacre — un ítem de cada bloc}

\begin{exemple}[com es resol un ítem ben fet, de cap a cap]
\textbf{Ítem tipus (Bloc 1):} els ingressos (€) de $5$ famílies són $\num{900}$,
$\num{1000}$, $\num{1100}$, $\num{1000}$, $\num{5000}$. Un \textbf{ítem complet}
respon: paràmetres \emph{i} interpretació.
\begin{itemize}[leftmargin=1.4em, itemsep=1pt]
  \item Mitjana: $\overline{x}=\dfrac{\num{9000}}{5}=\num{1800}\eur$.
  \item Mediana: ordenat $\num{900},\num{1000},\underline{\num{1000}},\num{1100},\num{5000}$,
        $\text{Me}=\num{1000}\eur$.
  \item \textbf{Interpretació:} el valor $\num{5000}\eur$ estira la mitjana; la
        \textbf{mediana ($\num{1000}\eur$) representa millor} el grup. \emph{Aquesta
        frase final és la que distingeix un AE d'un AS.}
\end{itemize}
\end{exemple}

\begin{exercicis}
\textbf{Resol el simulacre amb temps controlat. Un ítem per bloc:}
\begin{exlist}
  \item \textbf{(Bloc 1 — C1)} Per a les edats $18$, $20$, $19$, $21$, $52$: calcula la
        mitjana, la mediana i la desviació típica, i digues quina mesura de
        centralització representa millor el grup i per què.
  \item \textbf{(Bloc 2 — C2)} Un gràfic d'una taxa d'ocupació té l'eix vertical
        començant a $70$ i mostra una pujada de $72$ a $76$. Algú diu que «l'ocupació
        s'ha disparat». Quants punts ha pujat realment? Per què el gràfic exagera?
  \item \textbf{(Bloc 3 — C3)} Indica si la pregunta «\emph{No creus que caldrien més
        ajudes?}» és neutra o esbiaixada, reescriu-la si cal, i tabula aquestes $8$
        respostes (Sí/No): Sí, No, Sí, Sí, No, Sí, Sí, No.
\end{exlist}
\end{exercicis}

\subsection{Formulari-resum (suport)}

\begin{idea}
\textbf{Tingues-ho a mà mentre prepares la prova:}
\begin{itemize}[leftmargin=1.4em, itemsep=2pt]
  \item \textbf{Mitjana:} $\overline{x}=\dfrac{\sum x_i}{N}$ (suma dividida pel total).
  \item \textbf{Mediana:} \emph{ordena} i agafa el central (o la mitjana dels dos
        centrals si $N$ és parell).
  \item \textbf{Moda:} el valor que més es repeteix.
  \item \textbf{Desviació típica:} $\sigma=\sqrt{\dfrac{\sum (x_i-\overline{x})^2}{N}}$
        (no oblidis l'arrel!).
  \item \textbf{Lectura crítica:} mira la font, si l'eix comença a zero, «percentatge
        de què» i si es comparen grups de la mateixa mida.
  \item \textbf{Enquesta:} pregunta tancada, clara, neutra i d'una sola idea.
\end{itemize}
\end{idea}

\noindent\textbf{Pla d'estudi personal.} A partir de la teva graella d'autoavaluació
(sessió 5) i de com t'ha anat el simulacre, anota els \textbf{dos o tres tipus
d'exercici} que has de repassar amb més atenció abans de la prova:
\begin{center}
\renewcommand{\arraystretch}{1.5}
\begin{tabular}{c p{10.5cm}}
\toprule
\textbf{\#} & \textbf{Què he de repassar (i com)} \\
\midrule
1 & \rule{9.8cm}{0.4pt} \\
2 & \rule{9.8cm}{0.4pt} \\
3 & \rule{9.8cm}{0.4pt} \\
\bottomrule
\end{tabular}
\end{center}

% ---------------------------------------------------------------------
\fullsolucions{Sessió 15 — simulacre}
\begin{solucions}
\begin{sollist}
  \item \textbf{(Bloc 1)} Mitjana:
        $\overline{x}=\dfrac{18+20+19+21+52}{5}=\dfrac{130}{5}=26$. Mediana: ordenat
        $18,19,\underline{20},21,52$, $\text{Me}=20$. Desviació típica: distàncies al
        quadrat respecte de $26$: $64,36,49,25,676$, suma $850$;
        $\sigma^2=\dfrac{850}{5}=170$; $\sigma=\sqrt{170}\approx 13,04$. La
        \textbf{mediana ($20$) representa millor} el grup: el valor $52$ estira la
        mitjana fins a $26$ i infla molt la desviació típica.
  \item \textbf{(Bloc 2)} L'ocupació ha pujat $76-72=4$ punts. El gràfic exagera perquè
        l'eix vertical \textbf{comença a $70$} (no a $0$): així una pujada de $4$ punts
        ocupa molta alçada i \emph{sembla} enorme. Amb l'eix des de $0$, la mateixa
        pujada es veuria molt més moderada.
  \item \textbf{(Bloc 3)} La pregunta és \textbf{esbiaixada} («No creus que\dots»
        indueix el sí). Neutra: «Com valores el nivell actual d'ajudes?»
        (Insuficient / Adequat / Excessiu). Tabulació de les $8$ respostes:
        \textbf{Sí $=5$} ($62,5\%$), \textbf{No $=3$} ($37,5\%$).
\end{sollist}
\end{solucions}
